% Cote 119: offprints, reports, courses, CV and letters (typed and printed part): archive cote 119, site page 183 (PDF page 184).
% typescript. Draft compared: none.
% AI-transcribed by Claude Opus 5.5 (claude-opus-5-5, adaptive thinking, effort high), Message Batches API, batch r119_full (msgbatch_01DwZNm5Y7YsKBWyAN4sFs8w), 2026-11-13; not yet checked by eye.
\scanpage{119/183}
\mspage{- 2 -}

(5.5.1) \quad $f(z) = \sum_{n \geq 0} a_n \, (z-z_o)^n$ \quad où les coefficients sont donnés par
\[ a_n = \frac{1}{2 \pi i} \int_{|t-z_o|=r} \frac{f(t)}{(t-z_o)^{n+1}} \, dt \]
Faisons le changement de variable \quad $(t-z_o) = r e^{i\theta}$
\[ (7.1.1) \qquad a_n r^n = \frac{1}{2\pi} \int_0^{2\pi} e^{-n^{i\theta}} \, f(z_o + re^{i\theta}) \, d\theta \]

formule que l'on obtiendrait tout aussi bien à partir de
\[ f(z_o + re^{i\theta}) = \sum_{p \geq 0} a_p r^n e^{pi\theta} \]
en multipliant par $e^{-ni\theta}$ et en intégrant sur $[0,2\pi]$ car pour $r$ fixé la série
précédente est normalement \unsure{convergente}.

Si nous posons
\[ (7.1.2) \qquad M_{z_o}(r) = \sup_{0 \leq \theta \leq 2\pi} |f(z_o + re^{i\theta})| \quad \text{on obtient l'inégalité} \]
\[ (7.1.3) \qquad |a_n| \leq \frac{M_{z_o}(r)}{r^n} \quad \text{en majorant (7.1.1)} \]

Ce sont les \uline{inégalités de Cauchy}. On savait déjà (Séries Entières
début) que, comme $r < \rho$, $|a_n| r^n$ était majoré mais on a précisé ici la valeur
d'un majorant en fonction de $r$.

7.2. \uline{Théorème de Liouville}.
\uline{Théorème}. Une fonction $f(z)$ \unsure{holomorphe|hdomorphe} dans tout le plan \unsure{etbornée|et bornée} est constante.
\uline{Démonstration}: On applique (7.1.3) pour $z_o = 0$.
Si $|f(z)| \leq M$. en faisant tendre $r$ vers l'infini on voit que
$a_n = 0$ pour $n \geq 1$.
\uline{Remarque :} Si $P(z)$ est un polynôme non constant sans racine $\frac{1}{P(z)}$ est une
fonction holomorphe sur tout le plan dont on montrera simplement qu'elle est
bornée d'où une contradiction par le théorème de Liouville. On trouve donc une
seconde démonstration au théorème de d'Alembert- Gauss et le lecteur désireux
de faire la pige au jeune Gauss n'a plus que quelques autres démonstrations à

\margin{183}
